\documentclass{amsart}
% For dealing with references we use the comment environment
\usepackage{verbatim}
\newenvironment{reference}{\comment}{\endcomment}
%\newenvironment{reference}{}{}
\newenvironment{slogan}{\comment}{\endcomment}
\newenvironment{history}{\comment}{\endcomment}

% For commutative diagrams we use Xy-pic
\usepackage[all]{xy}

% We use 2cell for 2-commutative diagrams.
\xyoption{2cell}
\UseAllTwocells

% We use multicol for the list of chapters between chapters
\usepackage{multicol}

% This is generally recommended for better output
\usepackage{lmodern}
\usepackage[T1]{fontenc}

% For cross-file-references

% Package for hypertext links:
\usepackage{hyperref}

% For any local file, say "hello.tex" you want to link to please

% Theorem environments.
%
\theoremstyle{plain}
\newtheorem{theorem}[subsection]{Theorem}
\newtheorem{proposition}[subsection]{Proposition}
\newtheorem{lemma}[subsection]{Lemma}

\theoremstyle{definition}
\newtheorem{definition}[subsection]{Definition}
\newtheorem{example}[subsection]{Example}
\newtheorem{exercise}[subsection]{Exercise}
\newtheorem{situation}[subsection]{Situation}

\theoremstyle{remark}
\newtheorem{remark}[subsection]{Remark}
\newtheorem{remarks}[subsection]{Remarks}

\numberwithin{equation}{subsection}

% Macros
%
\def\lim{\mathop{\mathrm{lim}}\nolimits}
\def\colim{\mathop{\mathrm{colim}}\nolimits}
\def\Spec{\mathop{\mathrm{Spec}}}
\def\Hom{\mathop{\mathrm{Hom}}\nolimits}
\def\Ext{\mathop{\mathrm{Ext}}\nolimits}
\def\SheafHom{\mathop{\mathcal{H}\!\mathit{om}}\nolimits}
\def\SheafExt{\mathop{\mathcal{E}\!\mathit{xt}}\nolimits}
\def\Sch{\mathit{Sch}}
\def\Mor{\mathop{\mathrm{Mor}}\nolimits}
\def\Ob{\mathop{\mathrm{Ob}}\nolimits}
\def\Sh{\mathop{\mathit{Sh}}\nolimits}
\def\NL{\mathop{N\!L}\nolimits}
\def\CH{\mathop{\mathrm{CH}}\nolimits}
\def\proetale{{pro\text{-}\acute{e}tale}}
\def\etale{{\acute{e}tale}}
\def\QCoh{\mathit{QCoh}}
\def\Ker{\mathop{\mathrm{Ker}}}
\def\Im{\mathop{\mathrm{Im}}}
\def\Coker{\mathop{\mathrm{Coker}}}
\def\Coim{\mathop{\mathrm{Coim}}}

% Boxtimes
%
\DeclareMathSymbol{\boxtimes}{\mathbin}{AMSa}{"02}

%
% Macros for moduli stacks/spaces
%
\def\QCohstack{\mathcal{QC}\!\mathit{oh}}
\def\Cohstack{\mathcal{C}\!\mathit{oh}}
\def\Spacesstack{\mathcal{S}\!\mathit{paces}}
\def\Quotfunctor{\mathrm{Quot}}
\def\Hilbfunctor{\mathrm{Hilb}}
\def\Curvesstack{\mathcal{C}\!\mathit{urves}}
\def\Polarizedstack{\mathcal{P}\!\mathit{olarized}}
\def\Complexesstack{\mathcal{C}\!\mathit{omplexes}}
% \Pic is the operator that assigns to X its picard group, usage \Pic(X)
% \Picardstack_{X/B} denotes the Picard stack of X over B
% \Picardfunctor_{X/B} denotes the Picard functor of X over B
\def\Pic{\mathop{\mathrm{Pic}}\nolimits}
\def\Picardstack{\mathcal{P}\!\mathit{ic}}
\def\Picardfunctor{\mathrm{Pic}}
\def\Deformationcategory{\mathcal{D}\!\mathit{ef}}

\setlength{\emergencystretch}{3em}
\begin{document}
\title[Stacks Project, Tag 0GN5]{365 --- Flat coefficient base change and projection for higher quasi-coherent direct images}
\maketitle
\noindent Copyright (C) 2005--2025 Johan de Jong. Stacks Project authors. Source: \href{https://stacks.math.columbia.edu/tag/0GN5}{Tag 0GN5}.

\noindent Editorial difficulty note (not author text). Difficulty H2: The comparison must be constructed globally from derived adjunction and cup products even though the base-change morphism itself is not assumed flat. A finite-in-each-degree affine hypercovering is then needed to prove the isomorphism using flatness of the coefficient module over the original base..

\noindent Editorial provenance note (not author text). History: excerpt from revision \texttt{a0444\allowbreak 6e57e\allowbreak c1fbc\allowbreak 252a8\allowbreak 71afc\allowbreak ec775\allowbreak 2fb28\allowbreak 07b14}, coherent.tex, lines 1431--1549. Reference presentation was adapted; original-author context and core support blocks were retained or added. The mathematical wording is unchanged. Licensed under GNU FDL 1.2 or later; no invariant sections or cover texts. See ../licenses/COPYING. Context blocks reproduce author definitions and supporting results; they are not additional counted problems.

\begin{lemma}
\label{lemma-base-change-tensor-with-flat}
Consider a cartesian diagram of schemes
$$
\xymatrix{
X' \ar[d]_{f'} \ar[r]_{g'} & X \ar[d]^f \\
S' \ar[r]^g & S
}
$$
Let $\mathcal{F}$ be a quasi-coherent $\mathcal{O}_X$-module.
Let $\mathcal{G}$ be a quasi-coherent $\mathcal{O}_{S'}$-module
flat over $S$. Assume $f$ is quasi-compact and quasi-separated.
For any $i \geq 0$ there is an identification
$$
\mathcal{G} \otimes_{\mathcal{O}_{S'}} g^*R^if_*\mathcal{F} =
R^if'_*\left((f')^*\mathcal{G}
\otimes_{\mathcal{O}_{X'}} (g')^*\mathcal{F}\right)
$$
\end{lemma}

\begin{proof}
Let us construct a map from left to right. First, we have the base change map
$Lg^*Rf_*\mathcal{F} \to Rf'_*L(g')^*\mathcal{F}$. There is also the adjunction
map $\mathcal{G} \to Rf'_*L(f')^*\mathcal{G}$. Using the relative cup product
We obtain
\begin{align*}
\mathcal{G}
\otimes_{\mathcal{O}_{S'}}^\mathbf{L}
Lg^*Rf_*\mathcal{F}
& \to
Rf'_*L(f')^*\mathcal{G}
\otimes_{\mathcal{O}_{S'}}^\mathbf{L}
Rf'_*L(g')^*\mathcal{F} \\
& \to
Rf'_*\left(L(f')^*\mathcal{G}
\otimes_{\mathcal{O}_{X'}}^\mathbf{L}
L(g')^*\mathcal{F}\right) \\
& \to
Rf'_*\left((f')^*\mathcal{G}
\otimes_{\mathcal{O}_{X'}} (g')^*\mathcal{F}\right)
\end{align*}
where for the middle arrow we used the relative cup product, see
Cohomology, Remark \href{https://stacks.math.columbia.edu/tag/0B68}{remark cup product (Tag 0B68)}.
The source of the composition is
$$
\mathcal{G} \otimes_{\mathcal{O}_{S'}}^\mathbf{L} Lg^*Rf_*\mathcal{F} =
\mathcal{G} \otimes_{g^{-1}\mathcal{O}_S}^\mathbf{L} g^{-1}Rf_*\mathcal{F}
$$
by Cohomology, Lemma \href{https://stacks.math.columbia.edu/tag/08DE}{lemma variant derived pullback (Tag 08DE)}.
Since $\mathcal{G}$ is flat as a sheaf of $g^{-1}\mathcal{O}_S$-modules
and since $g^{-1}$ is an exact functor, this is a complex whose
$i$th cohomology sheaf is
$\mathcal{G} \otimes_{g^{-1}\mathcal{O}_S} g^{-1}R^if_*\mathcal{F} =
\mathcal{G} \otimes_{\mathcal{O}_{S'}} g^*R^if_*\mathcal{F}$.
In this way we obtain global maps from left to right in the
equality of the lemma. To show this map is an isomorphism we may
work locally on $S'$. Thus we may and do assume that $S$ and $S'$
are affine schemes.

\medskip\noindent
Proof in case $S$ and $S'$ are affine. Say $S = \Spec(A)$ and $S' = \Spec(B)$
and say $\mathcal{G}$ corresponds to the $B$-module $N$ which is assumed to be
$A$-flat. Since $S$ is affine, $X$ is quasi-compact and quasi-separated.
We will use a hypercovering argument to finish the proof; if $X$ is separated
or has affine diagonal, then you can use a {\v C}ech covering.
Let $\mathcal{B}$ be the set of affine opens of $X$. By
Hypercoverings,
Lemma \href{https://stacks.math.columbia.edu/tag/01H7}{lemma quasi separated quasi compact hypercovering (Tag 01H7)}
we can find a hypercovering $K = (I, \{U_i\})$ of $X$ such that each
$I_n$ is finite and each $U_i$ is an affine open of $X$. By
Hypercoverings, Lemma \href{https://stacks.math.columbia.edu/tag/01GY}{lemma cech spectral sequence (Tag 01GY)}
there is a spectral sequence with $E_2$-page
$$
E_2^{p, q} = \check{H}^p(K, \underline{H}^q(\mathcal{F}))
$$
converging to $H^{p + q}(X, \mathcal{F})$. Since each $U_i$ is affine
and $\mathcal{F}$ is quasi-coherent the value of
$\underline{H}^q(\mathcal{F})$ is zero on $U_i$ for $q > 0$. Thus the
spectral sequence degenerates and we conclude that the cohomology
modules $H^q(X, \mathcal{F})$ are computed by
$$
\prod\nolimits_{i \in I_0} \mathcal{F}(U_i)
\to
\prod\nolimits_{i \in I_1} \mathcal{F}(U_i)
\to
\prod\nolimits_{i \in I_2} \mathcal{F}(U_i)
\to \ldots
$$
Next, note that the base change of our hypercovering to
$S'$ is a hypercovering of $X' = S' \times_S X$.
The schemes $S' \times_S U_i$ are affine too
and we have
$$
\left((f')^*\mathcal{G} \otimes_{\mathcal{O}_{S'}} (g')^*\mathcal{F}\right)
(S' \times_S U_i) =
N \otimes_A \mathcal{F}(U_i)
$$
In this way we conclude that the cohomology modules
$H^q(X', (f')^*\mathcal{G} \otimes_{\mathcal{O}_{S'}} (g')^*\mathcal{F})$
are computed by
$$
N \otimes_A \left(
\prod\nolimits_{i \in I_0} \mathcal{F}(U_i)
\to
\prod\nolimits_{i \in I_1} \mathcal{F}(U_i)
\to
\prod\nolimits_{i \in I_2} \mathcal{F}(U_i)
\to \ldots
\right)
$$
Since $N$ is flat over $A$, we conclude that
$$
H^q(X', (f')^*\mathcal{G} \otimes_{\mathcal{O}_{S'}} (g')^*\mathcal{F})
=
N \otimes_A H^q(X, \mathcal{F})
$$
Since this is the translation into algebra of the statement
we had to show the proof is complete.
\end{proof}
\end{document}
